\begin{scanpage}{0014}{x}
significantly exceed the wavelength, the additional currents occupy a comparatively small part of its surface.

Since the current excited by the plane wave on an ideally conducting surface is distributed uniformly over it (the absolute magnitude of its surface density is constant) then the vector $\mathbf j^0$ may be called the ``uniform'' part of the surface current. The additional current $\mathbf j^1$ which is caused by the curve of the body's surface we will henceforth call the ``nonuniform'' part of the current. In the physical optics approach, only the uniform part of the current is considered. Therefore, it is no wonder that in a number of cases it gives unsatisfactory results. For a more precise calculation, it is necessary to also take into account the nonuniform part of the current.

In this book, the results of the author relating to the approximation solution of diffraction problems are discussed and systematized. Essentially, these results were briefly discussed in a number of papers [7--14]. Roughly at the same time, the works of other authors devoted to similar problems appeared. We will discuss them in more detail (in \S 25) after the reader becomes accustomed to the concepts being used in diffraction problems of this type. For the present, let us only note that in these works, as a rule, other methods are used.

In the book, problems of the diffraction of plane electromagnetic waves by complex metal bodies, the surfaces of which have discontinuities (edges), are investigated. The dimensions of the bodies are assumed to be large in comparison with the wavelength, and their surface is assumed to be ideally conducting.

Obviously, if the edges are sufficiently far from one another, then the current flowing on a small element of the body's surface in the vicinity of its discontinuity may be approximately considered to be the same as on a corresponding infinite dihedral angle (a wedge). In fact, in Chapter I it is shown (see also [5] \S 20) that the nonuniform part of the current on a wedge has the character of an edge
\end{scanpage}
